% Appendix C, part 1 of 3 (continued in parts 2 and 3); label prefix: evi. No new macros.
\section{Evidence: finite computations}
\label{sec:evi}

Nothing in this appendix is a step of a proof: the proofs of the note are complete without it, and no proof refers to it. It records the finite computations that were made while the statements of the note were developed, as evidence that the statements hold at the orders tested, together with their ranges and their controls. The arithmetic is exact over $\bbQ$ or over $\bbQ(\zeta)$, $\zeta=\re^{2\pi\ri/5}$, or it is done modulo a prime written out as a number; a zero test modulo a prime is probabilistic. A control is a deliberately wrong input that must make a check fail; we list the controls that failed as required, since they show that the checks can detect an error (a positive control is a modified input that must leave a check intact; these are named as such). Conifold frames are labelled by $(\alpha,\beta)$ as in \cref{def:set-frames}, $\varpi_0=f_0+\alpha\varpi_1+\beta f_2$; the two frames used most are $(0,0)$ and $(1/3,-2/7)$. Besides the canonical normalisation $P_{11}(0,z)=1$ of the gauge (\cref{def:gau-can}), some computations use another normalisation, $P=\Pc(1+m_5z^5+\dots)$, defined by $p_{5n,0}(0)=0$ for $n\geq 1$ (\cref{sec:sgn-norm}); its scalars are given in \cref{sec:evi-univ}. Several independent implementations were written; where a statement was checked by more than one of them this is said. A check at one genus or one order says nothing about the next.

\subsection{Calibration and genus zero}
\label{sec:evi-calib}

\emph{Calibration of the graph sum.} The graph sum of \cref{def:set-graphsum} (for the numbers $N_{1,d}$ of genus $1$, the corresponding genus-$1$ formula of the $R$-matrix action, which the note does not use), with the R-matrix normalised by (QRR), was compared with direct virtual localisation \cite{MR1666787} on the moduli spaces of maps to $\bbP^4$ at $\lambda_i=\zeta^i$. Both give $N_{2,0}=F_2(0)=-19/9600$, $N_{1,1}=-305/6$, $N_{1,2}=-535775/12$, $N_{2,1}=2129/5760$ and $N_{2,2}=203401/144$; in genus $2$ and degree one the split of the localisation contributions by the genera of the contracted components matches the graph sum graph by graph (the vertex graph contributes $-24161/5760$, the graph with two vertices of genus one $2629/576$). The graph sum predicts $N_{1,3}=-1540037125/36$, $N_{1,4}=-1082838731875/24$, $N_{2,3}=1016639975/384$ and $N_{2,4}=181712650375/48$. Localisation gives $N_{3,0}=F_3(0)=1/224000$; the genus-$3$ value of the graph sum at $q=0$ was not computed. The genus-$0$ numbers do not depend on the weights, at four choices of them. From the vertex class of \cite[(26)]{LP20} at $\lambda_i=\zeta^i$, exactly over $\bbQ(\zeta)$, one obtains again $N_{2,0}=-19/9600$ and $N_{3,0}=1/224000$, whereas the twist weight $+5$ in place of $-5$ gives $-97/28800$ and $1/81000$: this supports the sign of the twist weight in (QRR) independently of the graph sum (\cref{rmk:set-qrrsign}). The formal quintic is not the quintic threefold: the constant-map term of the latter at genus $2$ is $-5/144$.

\emph{Genus zero.} The mirror theorem (M) in the form of \cref{thm:set-mirror} was compared at order $q$ with degree-one localisation on $\overline{\mathcal M}_{0,1}(\bbP^4,1)$, as rational functions of $z$, at three random rational weight vectors and at $\lambda_i=\zeta^i$; the literal reading of \cite[Cor.~11.7]{MR1408320} fails by $+600$, and two further controls fail. The coefficients $\kappa^{\mathrm{LR}}_a$ of \cref{prop:set-cyclicLR} were computed to order $q^{30}$ by two implementations (one of them keeping the formal $\logq$ explicit, with fifteen periods); they are free of $\logq$ and have the stated closed forms, and the closing relation and the operator identity relating them to $\PFfour$ hold to that order. Four controls fail: the non-equivariant series, the twist $\cO(4)$ or $\cO(6)$ in place of $\cO(5)$, a wrong $\kappa^{\mathrm{LR}}_1$, and the guess $\kappa^{\mathrm{LR}}_3=C_1/C_0$. Degree-one localisation on $\overline{\mathcal M}_{0,2}(\bbP^4,1)$ at $\lambda_i=\zeta^i$, independent of (M), gives $\langle H,H\rangle_{0,2,1}=2875=5\coef{q}\kappa^{\mathrm{LR}}_1$ and $\langle H^3,H^4\rangle_{0,2,1}=-3250=5\coef{q}\kappa^{\mathrm{LR}}_3$, while the $\bbZ/5$-grading forces $\langle H^2,H^3\rangle_{0,2,1}=\langle H^4,H^4\rangle_{0,2,1}=0$ at $\zeta$ (the first is non-zero at generic weights). The Yukawa coupling $5L^5/(C_0^2C_1^3)$ reproduces the genus-$0$ invariants of the quintic threefold in degrees at most $5$ (instanton numbers $2875$, $609250$, $317206375$, $242467530000$, $229305888887625$).

\subsection{Graph sums of genus 2 to 5}
\label{sec:evi-graphs}

The outer terms $T_\Gamma$ of \cref{def:set-outer}, with $M=\Mo(\tau(u)\zh)$ truncated in $\zh$ at the order the genus needs (modulo $\zh^7$ and $u^4$ at genus $2$ and $3$, which suffices there; exactly in $u$, modulo $\zh^{3g-2}$, at genus $4$ and $5$), were computed for every leg-free stable graph of genus $2$, $3$, $4$ and $5$, modulo the primes $2^{89}-1$ and $2^{61}-1$ with rational reconstruction (some checks also modulo $2^{521}-1$), in the frames $(0,0)$ and $(1/3,-2/7)$. A term of genus $g$ is written $T_\Gamma=u^{2-2g}G_\Gamma(u)$ (\cref{lem:flh-hat}); an edge graph is regular if and only if $G_\Gamma=O(u^{2g-2})$, so only $G_\Gamma$ modulo $u^{2g-2}$ is needed. At genus $4$ and $5$ the outer sum was evaluated by splitting the edge of the R-matrix $R=R'M$ of the plain graph sum into the edge of $M$ and the edge of $R'$ conjugated by $M^{-1}$, $E_R(\psi,\psi')=E_M(\psi,\psi')+M^{-1}(\psi)E_{R'}(\psi,\psi')M^{-1}(\psi')^T$ (an identity because $M$ and, by \cref{thm:gam-D}, $R'$ are symplectic; \cref{def:set-outer}), the leaves of $R$ being unchanged: the plain sum over the graphs whose edges are marked accordingly, with Witten--Kontsevich numbers of a single colour at the vertices, was regrouped, by contracting the edges marked by $M$, into the outer graph sum. That this regrouping reproduces the outer graph sum is an argument that was checked against genuine vertex tensors only at genus $2$ and $3$. The results are in \cref{tab:evi-g2,tab:evi-g3,tab:evi-g4,tab:evi-g5}: in every genus and in both frames every edge graph is regular, and the vertex graph has the polar part $c_g\tc^{2-2g}$ and no other polar term. At genus $4$ and $5$ the tables give the totals; since the total is the sum of the vertex graph and the edge graphs, the statement on the vertex graph is there the same check as the regularity of the edge graphs, not a second one. At genus $2$ the total is $F_2^{\mathrm{cf}}=(1/1200)\tc^{-2}+0\cdot\tc^{-1}+31/30000+O(\tc)$ in the frame $(0,0)$; in the frames $(0,0)$, $(1,0)$, $(0,1)$, $(2/3,-1/7)$ the polar part is the same and the constant terms are $31/30000$, $31/30000$, $2749/120000$, $-251/120000$.

\begin{table}[ht]
\centering\small\setlength{\tabcolsep}{4pt}
\begin{tabular}{l|r|r|r}
graph of genus $2$ & plain, $\coef{\tc^{-2}}$ & plain, $\coef{\tc^{-1}}$ & natural, at $u=0$\\ \hline
one vertex of genus $2$ & $1807/3072000$ & $-791/15360000$ & $\tfrac{1}{1200}\tc^{-2}+O(1)$\\
genus-$1$ vertex with a loop & $-811/256000$ & $-1771/3840000$ & $479/756000$\\
genus-$0$ vertex with two loops & $-57/64000$ & $41/12800$ & $1/2250$\\
two genus-$1$ vertices, one edge & $843/1024000$ & $743/15360000$ & $-127/2016000$\\
genus $0$ with a loop, edge to genus $1$ & $3/51200$ & $-41/153600$ & $0$\\
two genus-$0$ vertices, three edges & $87/32000$ & $37/160000$ & $0$\\
two genus-$0$ vertices with loops, an edge & $9/12800$ & $-173/64000$ & $0$\\ \hline
sum & $1/1200$ & $0$ & \\
\end{tabular}
\caption{Genus $2$, per stable graph, summed over colourings, in the frame $(0,0)$, modulo both primes: the polar coefficients of the plain graph sum of \cref{def:set-graphsum} (vertices the TFT $\TFT$), and the terms $T_\Gamma$ of the natural form at $u=0$ (in the frame $(1/3,-2/7)$, modulo $2^{89}-1$, the edge graphs are regular as well). The $115$ coloured graphs give identical tables at both primes.}
\label{tab:evi-g2}
\end{table}

\begin{table}[ht]
\centering\small
\begin{tabular}{l|p{0.68\textwidth}}
stable graphs & $42$; with $0,1,\dots,6$ edges: $1,2,5,9,12,8,5$\\
edge graphs & all $41$ regular: frame $(0,0)$ modulo both primes, frame $(1/3,-2/7)$ modulo $2^{89}-1$\\
vertex graph, $(0,0)$ & $u^{-4}\bigl(\tfrac{1}{25200}-\tfrac{u}{9000}+\tfrac{407u^2}{3780000}-\tfrac{151u^3}{3780000}+\dots\bigr)=\tfrac{1}{25200}\tc^{-4}+O(1)$\\
vertex graph, $(1/3,-2/7)$ & $u^{-4}\bigl(\tfrac{1}{25200}-\tfrac{11u}{189000}+\dots\bigr)$, with the same polar part\\
\end{tabular}
\caption{Genus $3$, the natural form.}
\label{tab:evi-g3}
\end{table}

\begin{table}[ht]
\centering\small
\begin{tabular}{l|p{0.72\textwidth}}
stable graphs & $379$; with $0,1,\dots,9$ edges: $1,3,7,21,43,75,89,81,42,17$\\
edge graphs & all $378$ regular, in both frames, modulo both primes\\
total, $(0,0)$ & $u^{-6}\bigl(\tfrac{1}{180000}-\tfrac{7u}{300000}+\tfrac{701u^2}{18000000}-\tfrac{1459u^3}{45000000}+\tfrac{923639u^4}{67500000000}-\tfrac{171583u^5}{67500000000}\bigr)$ $=\tfrac{1}{180000}\tc^{-6}+O(1)$\\
total, $(1/3,-2/7)$ & $u^{-6}\bigl(\tfrac{1}{180000}-\tfrac{11u}{900000}-\tfrac{79u^2}{378000000}+\dots\bigr)$, with the same polar part\\
\end{tabular}
\caption{Genus $4$, the natural form; $1/180000=c_4$.}
\label{tab:evi-g4}
\end{table}

\begin{table}[ht]
\centering\small
\begin{tabular}{l|p{0.72\textwidth}}
stable graphs & $4555$; with $0,1,\dots,12$ edges: $1,3,11,34,100,239,492,784,1002,926,632,260,71$\\
edge graphs & all $4554$ regular: frame $(0,0)$ modulo both primes, frame $(1/3,-2/7)$ modulo $2^{89}-1$\\
total, $(0,0)$ & $u^{-8}\bigl(\tfrac{1}{660000}-\tfrac{7u}{825000}+\tfrac{199u^2}{9900000}-\tfrac{717u^3}{27500000}+\tfrac{449357u^4}{22500000000}-\tfrac{5593217u^5}{618750000000}$ $+\tfrac{14686042363u^6}{6496875000000000}-\tfrac{8370858511u^7}{32484375000000000}\bigr)=\tfrac{1}{660000}\tc^{-8}+O(1)$, by the Chinese remainder theorem from both primes\\
total, $(1/3,-2/7)$ & $u^{-8}\bigl(\tfrac{1}{660000}-\tfrac{u}{225000}+\tfrac{323u^2}{207900000}+\tfrac{66841u^3}{15592500000}-\dots\bigr)$, modulo both primes: different from the first at every order $u^1,\dots,u^7$, with the same polar part\\
\end{tabular}
\caption{Genus $5$, the natural form; $1/660000=c_5$.}
\label{tab:evi-g5}
\end{table}

\emph{Content.} At genus $4$, the $229$ graphs with at least six edges are regular a priori, because every outer edge and every outer leaf is $O(u)$ (\cref{rmk:mai-apriori}); of the other $149$ edge graphs, $81$ receive only contributions that vanish identically and $68$ vanish by cancellation among $6767$ non-zero contributions. At genus $5$, $2891$ graphs with at least eight edges are regular a priori, $1036$ receive no non-zero contribution and $627$ vanish by cancellation. Of the $3142$ polar coefficients at stake at genus $5$, $1655$ lie in graphs that receive contributions; of these, $517$ receive no non-zero contribution at that order and $1138$ vanish by genuine cancellation. A cancelling graph is a sum of between $31$ and $4525$ non-zero contributions (median $347$; the three one-edge graphs are sums of $3479$, $2296$ and $4525$).

\emph{Controls.} With the bare symbol $\Mo(\zh)$, that is $\tau=1$, all $149$ edge graphs of genus $4$ with at most five edges are singular. With $\tau$ replaced by $\tau+u^5$, exactly the three one-edge graphs of genus $4$ fail, at order $u^5$ of $G_\Gamma$, and the vertex graph acquires a term $\tc^{-1}$. At genus $5$ the controls were computed on the graphs with at most two edges: with the bare symbol all $14$ edge graphs among them are singular and the vertex graph acquires every polar order $\tc^{-7},\dots,\tc^{-1}$; the shift $\tau+u^7$ makes exactly the three one-edge graphs singular, at order $u^7$ only, and gives the vertex graph a term $\tc^{-1}$ only. A shift of $\tau$ by $cu^k$ moves a graph with $E$ edges at valuation at least $k+E-1$, by multilinearity in the edges; on a one-edge graph the lowest-order change is $c$ times a number that does not depend on $k$, and for the three one-edge graphs of genus $5$ these numbers, measured by the two controls ($k=1$ and $k=7$), are non-zero. So the computations at genus $5$ show that, within the ansatz $M=\Mo(\vartheta\zh)$ with $\vartheta\in\bbQ\fps{u}$, $\vartheta(0)=1$, the rescaling $\vartheta=\tau$ is the only one, modulo $u^8$, for which the one-edge graphs are regular (\cref{rmk:mai-tau}).

\emph{Independent implementations.} At genus $4$ a second implementation, written independently of the first, reproduced, modulo both primes in the frame $(0,0)$, all $379$ outer terms, all $379$ plain terms, every non-zero contribution and both controls (not the second frame); $26$ deliberate errors introduced into the first implementation were all detected, $23$ by the comparison with the values of genus $2$ and $3$ and the other three (in the coefficients of $u^4$ and $u^5$ of $\tau$ and in the orders $\zh\geq7$) by the checks at genus $4$. At genus $5$ an independently written implementation reproduced the contributions of $1474$ of the $4555$ graphs of the plain sum ($93551$ contributions modulo $2^{89}-1$, and $19591$ on $460$ of these graphs modulo $2^{61}-1$; no mismatch); the prime $2^{521}-1$ was used on $31$ of these graphs; and for each of the $4555$ graphs the contributions of its markings add up to its plain term, at both primes and in the second frame ($13665$ identities). The completeness of the graph list was checked by an independent induction, and the automorphism groups, the reconstruction, the Frobenius basis and $\tc$ were recomputed independently.

\emph{Caveats.} At genus $5$ the arithmetic is probabilistic (rational heights up to $2^{136}$ against the product $2^{150}$ of the two primes; a failure probability of about $2^{-138}$), and the two primes test the same implementation; the regularity in the second frame rests on one prime at genus $3$ and $5$; the coefficient of $u^6$ in $\tau$ is tested by no control at any genus, nor are the orders $\zh^{10},\zh^{11},\zh^{12}$ of $\Mo$; the two controls discriminate on $14$ of the $1663$ graphs of genus $5$ that carry content; the regrouping identity is checked as stated above, and the second implementation uses the same regrouping, so that an error in it would be common to both; the graph list is shared by the implementations. Finally, the value $c_g$ and the gap of the total do not depend on $M$: the total is independent of the factorisation, and the leading coefficient (of $u^{2-2g}$) of the vertex graph equals that of the plain total for every $M$ whose symbol at $u=0$ is the model's, because the outer edges are $O(u)$. So $c_g$ and the gap of the total are not evidence for the factorisation; only the regularity of the edge graphs tests it.

% Appendix C, part 2 of 3: linear Hodge frames, the normalisation, gauges. Continued in part 3.

\subsection{Linear Hodge frames}
\label{sec:evi-hodge}

At genus $2$ every coloured stable graph with an edge is singular at the conifold point in every linear Hodge frame tried (the true Hodge class, a linear Hodge class with a common weight, the exponent $+1$, the weight $5$, a colour-dependent class, and the plain graph sum), with the maximal pole $\tc^{-2}$, also for the three graphs whose vertices all have genus $0$. Every decoration has valuation exactly $-10(g-1)$ in $s=u^{1/5}$: the pole comes from the growth of the flat amplitudes $A_n$, of degree $4n+1$ in $L$, not from the $\Delta_i$; a change of Hodge frame by a $q$-independent scalar moves a row coefficient only at relative order $s^{4k}$, so the coefficients of $s^{-10},\dots,s^{-7}$ are those of the plain sum, non-zero for all $115$ coloured graphs (after summing over colourings the coefficient of $s^{-10}$ is the column $\coef{\tc^{-2}}$ of \cref{tab:evi-g2}). No constant covariant scalar frame, no constant diagonal frame (ten constants) and no constant symplectic colour-mixing frame (twenty-five parameters) regularises the edge graphs of genus $2$, modulo both primes; the factorisation by the bare symbol $\Mo(\zh)$, that is with $\tau=1$, kills the leading poles but leaves simple poles in the one-edge graphs; a covariant scalar correction of order $u$, with coefficients $(113/75,-21/125)$, regularises genus $2$ but has no common solution with genus $3$. These computations motivate the factorisation through the model; they are not used.

\subsection{The normalisation and the identity \texorpdfstring{$(\Gamma)$}{(Gamma)}}
\label{sec:evi-normal}

\Cref{tab:evi-windows} lists the finite windows in which $\Dn=1$ and the identity \eqref{eq:gam-Gamma} were tested directly, independently of the proofs of \cref{thm:gam-D} and \cref{thm:gam-main}, and some of the identities that enter those proofs.

\begin{table}[ht]
\centering\small
\begin{tabular}{>{\raggedright\arraybackslash}p{0.4\textwidth}|>{\raggedright\arraybackslash}p{0.3\textwidth}|>{\raggedright\arraybackslash}p{0.22\textwidth}}
statement & range & arithmetic\\ \hline
$\Dn_0=1$, through the connection form of \cref{sec:con-finite} & $z^{204}$ ($g\leq 69$), two implementations; $z^{125}$, a third & exact, frame $(0,0)$\\
$\Dn_0=1$ & $z^{1004}$ ($g\leq 335$) & modulo $2^{89}-1$, $2^{61}-1$\\
the actual gauge is $\Pc$ (\cref{cor:mai-actual}) & $z^{35}$; $z^{150}$ & two primes, two frames\\
\eqref{eq:gam-Gamma}: $c_n(1)=h_n$ & $n\leq 60$ (Laurent expansion at $L=1$); $n\leq 40$ (tracked precision); $n\leq 30$ (fit at $32$ points) & exact over $\bbQ$\\
\eqref{eq:gam-Gamma} & $z^{500}$ & modulo $2^{127}-1$, $2^{107}-1$\\
the chart of \cref{thm:gam-lmi} at $L=1$ equals $h$ & $z^{60}$ & exact over $\bbQ(\zeta)$\\
$\PhiSP=\Phican$ (\cref{thm:con-Pb}) & $\zh^{80}$, $u^{75}$ & modulo $2^{61}-1$ (two frames), $2^{89}-1$\\
$\cV=\varpi_0(\YCQ)_0$ (\cref{prop:con-res} with (G)) & $u^{40}$; $u^{30}$ & exact, two frames; two primes\\
$c(z;0)=\Ginf$ (\cref{thm:gam-ginf}) & $z^{60}$ (chart at $L=0$); $z^{80}$ (the value) & exact\\
$\fc^{\mathrm{can}}_m=\coef{z^{5m}}\Ginf$ (\cref{cor:gam-free}) & $m\leq 25$; $m\leq 100$ & exact; two primes\\
\end{tabular}
\caption{Finite windows of $\Dn=1$ and of the identity \eqref{eq:gam-Gamma}.}
\label{tab:evi-windows}
\end{table}

\emph{Theorem~P and its ingredients.} \Cref{thm:con-Pa}~(i)--(iv) at $L_0=2,3,-2,1/2,5/3$, that is also at $u=32$ and $u=-1/32$, outside $0<u<1$: $c_n(L_0)=L_0^{4n}\mathrm{ph}^{\mathrm{SP}}_n(L_0^{-5})$ exactly for $n\leq 9$, with $c_n(L_0)$ computed by one-variable Wick contraction and $\mathrm{ph}^{\mathrm{SP}}_n$ from the recursion of \cref{lem:gau-HPF}; controls (no half-density, the wrong sign of the twist term, the scaling $L_0^{4n-1}$) fail at $n=1$. By a second algorithm (a Morse coordinate and Lagrange inversion), modulo $2^{61}-1$: the same identity at twelve random points for $n\leq 12$, the derivative identity of the proof of \cref{thm:con-Pa}~(i) for the amplitudes $\sigma^m$, $m\leq 4$, $n\leq 8$, at six points, and $\PFfive^{L}(L\,c)=0$ through $z^7$ at three points; controls (a shifted free coefficient, the scaling $L_0^{4n-1}$, the amplitude $q(L)$ in the derivative identity, no half-density) fail. \Cref{lem:con-S2} through $t^{60}$ modulo two primes and through $t^{80}$ exactly; the uniformising map of \cref{lem:con-cfu}, with $V=\tau^{-5/4}$, through $u^{20}$; exactly, the relabelling of \cref{lem:con-relab} on $965$ hatted coefficients, the identification of \cref{lem:con-id} on $4283$ coefficients and through $\zh^{25}u^{25}$ in its gauge-free form, and the key lemma (\cref{lem:con-key}) on $884$ keys.

\emph{The degree.} For \cref{prop:con-deg}: for $n\leq 30$ and two choices of the free coefficients, every resonance is trivial and every $A_n$ lies in $L\bbQ[L]$ with the right grade; with a lower amplitude containing a term of non-positive valuation the resonance is obstructed at $n=5$ and $n=10$, and the twist factors $(1,2,3,4,10)$ break the valuation property; a second implementation checked the operator identities and the recursion for $n\leq 40$. The bound $\deg\mathrm{ph}_n\leq\lfloor 4n/5\rfloor$ is an equality for every $n\neq 2$ with $n\leq 1004$ (two primes) and $n\leq 300$ (exactly), and $\deg\mathrm{ph}_2=0$; the equality is not claimed in the note.

\emph{Theorem~U and the connection form.} The construction along the proof of \cref{thm:gau-U} has every coefficient a power series and every resonance unobstructed through $z^{80}$ modulo two primes (and for $k\leq 100$, resp.\ $k\leq 40$, by a second implementation, modulo two primes, resp.\ exactly). Three formulas for the coefficient of $z^{5n}$ of the scalar $\cn(z)=P_{11}(0,z)$ of a gauge agree for $n\leq 12$, exactly, in two frames.

\emph{The natural factorisation.} For \cref{cor:mai-natfact}, from rows built with the free coefficients $\coef{z^{5m}}\Ginf$, in the frames $(0,0)$ and $(1/3,-2/7)$, through $u^8$ and $\zh^{11}$, modulo $2^{61}-1$: the rows close; $\Psi^T\eta^{-1}\Psi=1$; $G(u,0)=G(0,\zh)=1$; $\coef{u^n}G$ has no term beyond $\zh^{\lfloor(5n+4)/4\rfloor}$; $\tilde R$ and $G$ are symplectic. Placing $\Mo$ on the left, or omitting $\tau$, fails at $(u^1,\zh^3)$, $(u^2,\zh^4)$, $(u^3,\zh^5)$, and $\gamma=7/10+1/3$ breaks the closing of the rows. Shifting the first free coefficient by $1/7$, which makes $\Dn_0\neq 1$ from $z^5$ on, leaves (a), (b), the polynomiality of $G$ and $G(0,\zh)=1$ intact; the first failure is $G(-\zh)^TG(\zh)\neq 1$, at $(u^8,\zh^{10})$. So among the conclusions of \cref{cor:mai-natfact} the symplecticity of $G$ is the only one that fails for this perturbation, which breaks $\Dn_0(z)\Dn_0(-z)=1$; it is blind to the factors $\exp(a(z))$ with $a$ odd (\cref{rmk:con-why}). Earlier windows of the factorisation: $\zh\leq 24$, $u\leq 17$ (two primes, $215$ vanishing coefficients) and seven frames at $\zh\leq 12$, $u\leq 8$; in both, the computed degrees satisfy even the bound $\lfloor(5n+2)/4\rfloor$; the placement of $\Mo$ on the left, and the rescalings $\tau=1$, $(1-u)^{4/9}$, $(u/\tc)^{4/5}$, fail at the first order where they differ from $\tau(u)$.

\emph{Controls of the normalisation.} The other normalisation differs from the canonical one first at $\zh^5u^4$, by $m_5$, and changing the first free coefficient by $m_5$ is invisible to every regularity test; $\tau$ replaced by $1$ (that is $v=u$) breaks regularity, obstructs the uniformising map at its first order and leaves the gauge equation without solution; weights without the twist weight $-5$ break $\mathrm{ph}_1(1)=h_1$; replacing $B_8$ by $B_8+1$ in $h$ breaks $\mathrm{ph}_n(1)=h_n$ first at $n=7$, and replacing $B_{96}$ by $B_{96}+1$ leaves everything through $z^{94}$ intact and changes the coefficient of $z^{95}$ of $\Dn_0$ to $5^{-95}/9120$; dropping the Stirling series of $\Gamma(1+5\sigma/z)$, or reversing the signs of the $\bk_k$, fails at $\zh^1$, and dropping that of the factor of index $j=0$ fails at $z^1$.

\emph{Independence.} Further conifold frames give no independent evidence for $\Dn=1$, since the connection data do not depend on the frame (\cref{cor:con-frame}); regularity checks are automatic given Theorem~U; only the odd coefficients $a_j$ of $\log\Dn_0=\sum_ja_jz^{5j}$ would be independent if $\Dn_0(z)\Dn_0(-z)=1$ were known a priori, which it is not (\cref{rmk:con-genus}); and the direct and the reduced tests share the Picard--Fuchs operator, the constants of (QRR) and the model.

\subsection{Cyclic systems and gauges}
\label{sec:evi-gauge}

\emph{The cyclic form.} The factorisation of \cref{thm:gau-cyc} holds exactly over $\bbQ$ below $u^{57}$ in the frames $(0,0)$, $(1/3,-2/7)$, $(7/11,-3/5)$; each lemma of \cref{sec:gau-cyc} separately to $u^{70}$ in the frames $(0,0)$, $(1/3,-109/210)$, $(7/11,-23/22)$; below $u^{77}$ in five frames and identically in $(\alpha,\beta)$ below $u^{27}$, with $\det\Wr_{\mathbf f}=-2/u^2$ for the standard basis $\mathbf f$ of \cref{def:gau-conc}; in $17$ frames to $u^{74}$, among them power-series solutions $y_0$ with $y_0(0)=0$ in place of $\varpi_0$ (the scope of \cref{thm:gau-cyc}), these to $u^{68}$; and by two further implementations (to $u^{70}$ in ten frames, among them some with $y_0(0)=0$, and to $u^{60}$ in four frames). Controls fail: non-isotropic pairs, a wrong value of $\gamma$, and $\PFfour+u\Dq$, which is not self-adjoint. The closing relation holds on the rows expanded in $s$, in every colour, for $z^1,\dots,z^6$ and at least $32$ orders of $s$, modulo two primes.

\emph{Theorem~U and the gauge identity.} Exactly over $\bbQ$ for $k\leq 16$: every $p_{k,b}$ is a power series, resonances occur only at $k\equiv 0$ modulo $5$, and the residues are $(0,1/5,-3/5,-2/5,-1/5)$. Built along the proof of \cref{thm:gau-U} in both frames for $z^0,\dots,z^{15}$: the gauge equation in every entry, $P_{11}(0,z)=1$, $p_{0,a}(0)=1$ and $\prod_ap_{0,a}=1$ to $u^{30}$, and $p_{0,0}=\tau^{-1/4}/\varpi_0$ (\cref{lem:gau-diag}); the entry $(\Pc)_{30}(0,z)=\tfrac{127}{3150}z^3+\tfrac{1416472237}{725359556250000}z^8+\dots$ in the frame $(0,0)$ and $-\tfrac{239}{450}z^3-\tfrac{20846188763}{725359556250000}z^8+\dots$ in the frame $(1/3,-2/7)$, and $p_{1,b}(0)=(8/45,0,8/45,-1/15,-1/15)$, resp.\ $(-7/45,0,-7/45,-1/15,-1/15)$. An independent recursion in $u$ with unknown $P(0,z)$ ($1574$ linear equations over $\bbQ$) has, below $z^{24}$ and in both frames, exactly the solutions $\Pc\,\cn$ with $\cn$ a series in $z^5$. The uniqueness clause, including non-graded gauges, modulo $2^{61}-1$: the projection of the solution space onto a window has dimension $\#\{k\leq 15:k\equiv\delta\bmod 5\}$ in degree $\delta$ for the quintic and the pulled-back model, and, by a second implementation, $\#\{k\leq 10:k\equiv\delta\bmod 5\}$ for random one-pole systems; without the pole the solutions form a module of rank $5$; with unequal products the solution space is zero on the window tested. The identity (G) of \cref{thm:gau-G}, exactly over $\bbQ$, $z^{-38},\dots,z^{20}$, $u<29$, two frames: $176$ coefficients of positive powers of $z$ vanish ($152$ by genuine cancellation of up to $142$ terms) and $1038$ coefficients equal those of $\YCQ$; in the other normalisation $127$ of the $176$ fail, and $\Pc(1+z^5/7)$ violates (G) at $z^5$. The canonical gauge is symplectic for $z^0,\dots,z^{20}$ modulo $u^{29}$, while the other normalisation has the defect $(\cn(z)\cn(-z)-1)\eta^{-1}$, first at $z^{10}$. The vanishing solutions (\cref{lem:gau-vanmod}, \cref{lem:gau-vanq}) and the diagonal: valuations for $n\leq 45$, $p_{0,0}$ to $u^{37}$, closed forms for $m\leq 80$.

\emph{Model rows, symplecticity, colours.} The symbol of the rows equals the model in all rows and colours for $n\leq 6$ (two primes); the stationary phase equals $\Fs_0$ through $t^{10}$ exactly; all rows through $\zh^{18}$ and the identity row through $\zh^{40}$. $\Mo(-t)^T\Mo(t)=1$ through $t^{24}$ in $\bbQ(\zeta)$; controls fail: a diagonally rescaled, gauge-equivalent system that is not $\eta$-self-adjoint, the pairing without $t\mapsto -t$, and the pairing on $a+b\equiv 2$. The colour-$i$ T-system has the solution $\zeta^{ib}\Fs_b(\zeta^{-i}t)$, $\Fs_0=1-\tfrac{3}{20}t+\tfrac{9}{800}t^2+\tfrac{11571}{80000}t^3-\tfrac{2415609}{6400000}t^4+\dots$, $\Fs_b(0)=(1,-1,-1,1,1)$, to $t^{10}$ exactly in $\bbQ(\zeta)$ (\cref{lem:gau-colmod}). The formal WKB columns of $\PFfive$ at large radius, colour $i$, $z^0,\dots,z^6$, $q^0,\dots,q^8$, over $\bbQ(\zeta)$, with the constants fixed by $h_{\bw_i}$: no logarithm, and the covariance $\Phi_i(\zh,u)=\Phi_0(\zeta^{-i}\zh,u)$ of \cref{lem:gau-colrows} in every coefficient; controls (colour $i$ normalised by $h_{\bw_0}$; the operator $\PFfive+qz^2$) fail. The transport operator of colour $i$ in the proof of \cref{lem:gau-wkb} acts as $Lg\mapsto 5\zeta^{4i}\Dq g$; this was checked on the monomials $L^m$, symbolically in the colour, with the same control.

% Appendix C, part 3 of 3: the model, transport, universality, local Mellin inversion. Last part of appendix C.

\subsection{The model}
\label{sec:evi-model}

\Cref{tab:evi-model} lists the main checks of \cref{sec:mod}.

\begin{table}[ht]
\centering\small
\begin{tabular}{>{\raggedright\arraybackslash}p{0.4\textwidth}|>{\raggedright\arraybackslash}p{0.36\textwidth}|>{\raggedright\arraybackslash}p{0.16\textwidth}}
statement & range & arithmetic\\ \hline
topological recursion on $x=\Wv$ is $(-5)^{h-1}$ times the model CohFT, i.e.\ the model correlators are $(-5)^{1-h}$ times those of topological recursion (\cref{thm:mod-tr}) & $72120$ colour-basis correlators with $\psi$-deficit $\leq 6$ at $(h,n)=(0,3)$, $(0,4)$, $(0,5)$, $(1,1)$, $(1,2)$, $(1,3)$, $(2,1)$, $(2,2)$, $(3,1)$ ($4665$ non-zero); on $x=\kappa\Wv$, $\kappa=1,4,-1$: $6495$ each & two primes\\
the R-matrix of the curve $x=\kappa\Wv$ is $\Mo(z/(\kappa\varsigma^4))$ (\cref{prop:mod-Rtr}) & all $25$ entries through $z^{10}$, $\kappa=1,4,-1$, $\varsigma=1,2$ & exact, $\bbQ(\zeta)$\\
the dilaton-leaf sign of \cite[Def.~2.13]{PPZ15} & $(h,n)=(1,1)$, four points, $\kappa=1,4$; the opposite sign fails at all four & exact\\
$\FTR_h(v=1)=B_{2h}/(2h(2h-2))=-\tfrac{1}{240},\tfrac{1}{1008},-\tfrac{1}{1440},\tfrac{1}{1056}$ & $h=2,\dots,5$ & $\bbQ(\zeta)$ ($h\leq 3$); two primes\\
the period identity of \cref{prop:mod-period}; $\Sigma_h=-4^hB_{2h}/(2h)$ & $h\leq 7$; $h\leq 20$ & exact\\
the lattice theorem (\cref{thm:mod-lattice}) & weight $\leq 400$; $\leq 1000$; $16200$ raw-leg kernel vectors & exact; two primes\\
the leg dictionary (\cref{prop:mod-LD}) & $35000$ pairings per prime; $52$ frame one-point functions, $h=2,3$ & two primes; $\bbQ(\zeta)$\\
the Model Lemma on $\sK$ (\cref{thm:mod-MLK}) & $5289$ correlators with $h\leq 4$, including $1302$ at $(h,n)=(2,4)$; $1395$ singular, all zero & exact; two primes\\
\end{tabular}
\caption{Checks on the model.}
\label{tab:evi-model}
\end{table}

Further, modulo two primes unless stated: the limits $v\to 0$ of the correlators at $13$ pairs $(h,n)$, for instance $\omega_{1,1}\to-\tfrac58\dd p/p^6$ and $\omega_{2,1}\to\tfrac{2079}{128}\dd p/p^{16}$; the closed formulas of the correlators against the recursion, exactly at $(0,3)$, $(1,1)$, $(0,4)$, $(1,2)$, $(2,1)$; the formula of \cref{rmk:mod-hock} for $h\leq 12$ (exactly). The $28$ non-zero pairings of legs of the class $\cL$ at $(1,1)$, $(2,1)$, $(1,2)$, $(2,2)$ have $v$-exponent $\geq 0$, while legs outside $\cL$ are singular (for example the pairing of $\omega_{1,1}$ with $-\log p$ is $-\tfrac1{12}v^{-1}$). The residues $\Res_{\varsigma_b}x^k\omega_{1,1}$ vanish for $b=0,\dots,4$ and $k=0,\dots,6$, while $40$ of the $45$ residues of $p^k\omega_{1,1}$, $k=1,\dots,9$, do not, and $\Res_{p=0}x^5\omega_{1,1}=-15/4$, so that the circle pairing of $x^5$ is not zero (\cref{def:mod-pair}). In the base case of \cref{thm:mod-V} the local formula of \cref{lem:mod-w11} gives $-1/60$ at each branch point at $v=1$, in total $-1/12$ (the opposite orientation of the ray gives $+1/12$), and the integral of the closed form of $\omega_{1,1}$ gives $-1/(12v)$ at $v=1,32,-3$. No log-free leg lies outside $\cL$ among $10400$ (weight $\leq 320$) and $16200$ (weight $\leq 400$) vectors, exactly; \cref{prop:mod-dii} holds for $k\leq 100$ and $n\leq 2000$ exactly and $n\leq 3000$ modulo two primes; \cref{prop:mod-V1} holds on $32559$ tuples with $h\leq 4$ ($21733$ by a second implementation), and the wrong-sign controls fail. The singular lists of the Model Lemma at genus $0,1,2,3$ (at most $5,5,4,3$ insertions, frames with $k\leq 3h+3$, $\psi$-deficit $\leq 6$; $43$, $418$, $425$, $127$ lists) all vanish: $326$ by genuine cancellation, $642$ with every term zero although allowed, $45$ trivially; deliberately wrong frames give non-zero values. The $17416$ singular ordered tuples of legs in $\cL$ and $-\log p$ at $(4,2),(3,3),(2,4),(1,5),(0,6)$ vanish modulo each of the two primes, and so do $40739$ multisets with $h+n\leq 7$, $h\leq 6$, by a second implementation.

\emph{Off $\sK$} (\cref{rmk:mod-offK}), exactly at $v=1$: $\langle\psi^5H\rangle_3=53/1209600=-\tfrac{53}{192}(\Fm_3)'$, whereas the one-point law of \cref{cor:mod-onepoint} would give $\tfrac1{96}(\Fm_3)'=-1/604800$; at genus $4$ the remainders of $\psi^5H$, $\psi^6\,1$, $\psi^7H^4$ are $\tfrac{19}{100800000}(1,-1,-\tfrac12)$; and the quotient $(\langle XY\rangle_h-(\Fm_h)''\Pi_0(X)\Pi_0(Y))/(\Fm_h)'$ depends on $h$, for instance $-23/12$ at $h=2$ and $-13/48$ at $h=3$ for $X=\psi H$, $Y=\psi^2H^3$.

\subsection{Transport and the comparison}
\label{sec:evi-transport}

\Cref{lem:tra-E} holds exactly over $\bbQ$ for random symplectic $R=\exp(X)$, generic and $\bbZ/5$-graded, with random $f$ in both slots: the numerator is divisible by $\psi+\psi'$ if and only if (Sp) holds, and the two sides agree; two further implementations, with $R=\exp(X)$ and with $R$ the Cayley transform of a random $X$, modulo $u^8$ and $u^7$, for $f$ built from $\YC$ and for random $f$ with the required decay, agree on all coefficients. Controls fail: a sign flip, a non-symplectic $R$, and an $f$ without decay, which gives different values for different truncations. For $\Pc$, in both frames, $z^0,\dots,z^{16}$, $u<16$, exactly: (Sp); (G) on $575$ complete coefficients; $T_P=\tc$ through $u^{14}$; $\Pi_0(q_T)=\tc$ and $\Pi_1(q_T)=\dots=\Pi_7(q_T)=0$; the edges divisible by $\psi+\psi'$ and annihilated by $\Pi_0,\dots,\Pi_7$ on $2985$ coefficients (\cref{cor:tra-edges}, \cref{lem:tra-T}); in the other normalisation (G) fails on $77$ coefficients, $\Pi_5(q_T)\neq 0$, the edge condition fails on $92$, and symplecticity fails from $z^{10}$. For the outer rows $R'$ of the natural form ($\zh\leq 24$, $u<18$, both primes, both frames) the edges are annihilated by $\Pi_k$ for $k=0,\dots,20$ ($2160$ coefficients each) and the leaf point by $\Pi_1,\dots,\Pi_{20}$ modulo $u^{18}$; with the scalar $\cn_5=1/7$ every condition fails and the leaf point acquires exactly a $\Pi_5$ component. In both frames $\delta_P=dT_P-v=u^2/30+(127/3150)u^3+\dots$ for $P=\Pc$ (it does not depend on the frame: $T_{\Pc}=\tc$ by \cref{prop:tra-can} and $d=1/p_{0,0}=\tau^{1/4}\varpi_0$ by \cref{lem:gau-diag}, so $dT_{\Pc}=\tau^{1/4}\varpi_1$); hence $\tc-\hat v=\delta_P/d=u^2/30+O(u^3)$, $\hat v=v/d$, the value used in \cref{rmk:red-leaf}.

\emph{Frames.} $\YC$ from the closed form of \cref{lem:gau-vanshape} agrees with the regular solution of the model system ($n<96$); the frames with $k\leq 80$, built with $\Theta^{-1}=\sum_r(-N)^r$ (\cref{lem:tra-theta}) and by elimination, agree ($324$ frames), satisfy $\Pi_i(\fr(k,a))=0$ for $0\leq i\leq k+12$, and have integral exponents $\geq 1$; by generic elimination at $v=1,2$ for $k\leq 60$ the trivial frames are exactly the ten with $k<\nu_a+1$ (\cref{lem:tra-frames}), and the rule $k\leq\nu_a+1$ is refuted. Controls fail: frames built from $\YC(+z)$ ($304$ of $314$ not annihilated) and $\YC$ without its terms $z^{-n}$, $n\geq 5$ ($214$ of $224$). In the $u$-adic setting, modulo $u^{10}$ (\cref{lem:tra-conv}): $\Theta$ is invertible, frame decompositions reconstruct random elements, $3312$ orders of frame coefficients equal $(4(k-j)+\nu_a+1)/5$, and a test bilinear form converges with the sharp bound $4K/5-3$; with $v=1+u$ in place of $v(u)$ there is no $u$-adic limit.

\emph{Theorem~R and the comparison.} In the flat sum of $\Pc$ (\cref{def:red-flatsum}) at genus $2$, frame $(0,0)$, the vertex graph is $u^{-2}(\tfrac1{1200}-\tfrac{7u}{6000}+\tfrac{287u^2}{864000}+\dots)$ and the six edge graphs are regular, with the values of \cref{tab:evi-g2}; in the frame $(1/3,-2/7)$ the same holds with the polar part $(1/1200)\tc^{-2}$. Controls: removing the leaves gives $-\tfrac1{900}u^{-1}$; the dressings $(1+u^mzE_{21})P$, $m=0,1$, are singular from order $u^m$; the non-gauge dressings $\exp(u^mzN)P$, $N=E_{04}+E_{43}$, are regular, with $T_P=u+u^2/5+\dots$ for $m=0$ and $T_P=\tc+O(u^3)$ for $m=1$. The comparison theorem (\cref{thm:flh-main}) at genus $3$: the flat and the outer terms of the $42$ graphs agree on the coefficients of $u^{-4},\dots,u^3$ ($336$ per prime and frame; frame $(0,0)$ modulo both primes, frame $(1/3,-2/7)$ modulo $2^{89}-1$), of which $108$ coefficients of edge graphs and $4$ of the vertex graph are non-zero; a second implementation reproduced all $252$ non-zero coefficients of the regrouping. Controls change many coefficients: a constant $\Mo$ in the splitting ($94$ of $252$), the dilaton factor dropped ($77$ of $336$), the leaves removed ($88$), the sign of the edge dropped ($103$), the other normalisation ($21$, in six graphs), $R(-z)$ for $R(z)$ ($112$), the dressing $\exp(uzN)$ ($74$).

\subsection{Universality and necessity}
\label{sec:evi-univ}

\emph{The other normalisation} (\cref{rmk:mai-stored}, \cref{rmk:mai-graphlocal}). Its scalar series $\cn(z)=P_{11}(0,z)=1+m_5z^5+m_{10}z^{10}+\dots$ has, exactly over $\bbQ$ and in both frames, $m_5=-43397/807975000$ and $m_{10}=-765368951400317/1305647201250000000$, so that $\cn(z)\cn(-z)=1+(2m_{10}-m_5^2)z^{10}+\dots$ with $2m_{10}-m_5^2=-18223115111903/15543419062500000$. At genus $4$ exactly six edge graphs are singular, the two one-edge bridges (between vertices of genus $1$ and $3$, and of genus $2$ and $2$) from order $u^4$ of $G_\Gamma$, and at order $u^5$ the genus-$3$ vertex with a loop and three graphs with two edges, while the other $372$ terms are unchanged and the polar part of the total does not move. Per unit $m_5$ the responses at the orders $u^4$, $u^5$ are $-\tfrac{1}{8400000},-\tfrac{187}{680400000}$ for the vertex graph, $-\tfrac{1}{1008000},\tfrac{13}{6804000}$ and $\tfrac{1}{900000},-\tfrac{53}{32400000}$ for the two bridges, and $-\tfrac{41}{68040000},-\tfrac{11}{34020000},\tfrac{49}{16200000},-\tfrac{17}{8100000}$ for the four graphs at $u^5$; they sum to zero at both orders, and the responses at $u^4$ are the same for perturbed systems. Two implementations agree.

\emph{Universality} (\cref{cor:mai-univ}, \cref{rmk:mai-annih}, \cref{lem:mai-coord}). Canonical gauges of perturbed one-pole systems that do not come from a Picard--Fuchs operator make all $378$ edge graphs of genus $4$ regular, and the vertex graph is then $(1/180000)(T/\lambda)^{-6}$ plus a regular series; likewise six perturbed systems at genus $2$ and $3$, in the system's own coordinate $T$. Annihilating non-gauge dressings are regular at genus $2$ ($6$ graphs), $3$ ($41$) and $4$ ($378$), with the vertex graph $c_gT_{\mathrm{eff}}^{2-2g}$, for instance $T_{\mathrm{eff}}=u+\tfrac{49}{45}u^2+\dots$ against $\tc=u+\tfrac{7}{10}u^2+\dots$; dressings that violate annihilation or flat regularity, and random orbits of the first column, give singular edge graphs (for example $14$ at genus $3$). Changes of the product $\Bprod$: for eight systems all edge graphs are regular at genus $2$ and $3$ modulo both primes, and at genus $4$ (those with at most two edges, and all $378$ for two systems), with the vertex graph in the system's own coordinate; $\phi'^5\Bprod_X(\phi)=\Bprod_q$ holds exactly for the change of coordinate $\phi$ of \cref{lem:mai-coord}; the controls fail as predicted. At first order, at genus $2$ ($9$ variables) the dressings that keep every edge graph regular are exactly the annihilating ones; at genus $3$ ($30$ variables) annihilation implies regularity, and the two further regular directions break the vertex law and are singular at genus $4$; without flat regularity ($20$, resp.\ $68$, variables) annihilation no longer implies regularity ($5$, resp.\ $23$, independent responses survive).

\subsection{Local Mellin inversion and the formal lemmas}
\label{sec:evi-lmi}

\emph{Theorem LMI on general data} (\cref{thm:gam-lmi}), exactly over $\bbQ$: for $\varphi=0$, $\cF=1$ the series $c^{\cF}$ is $1$ ($M_*=-2,\dots,3$, through $z^{12}$); in ten random cases ($a_*\in\{-1,2,-1/3,5/2,5,-7/4\}$, $M_*\in\{-1,0,1,2\}$) and four cases with entire data, the one-variable series equals the chart series away from $s_*$, the chart at $s_*$ equals \eqref{eq:gam-lmi}, and the Laurent expansion at $s_*$ has no pole (through $z^{12}$ and $z^{14}$); twelve further random cases (among them the non-integer values $M_*=1/3,-2/5,7/2$, outside \cref{def:gam-lmidata}; cf.\ \cref{rmk:gam-Mcomplex}) and a case with $\varphi''(s_*)=0$, through $z^{10}$ exactly and $z^{12}$, $z^{14}$ modulo two primes, and four more cases modulo these primes, two of them with $a_*=\sqrt{-1}$; the two-variable chart equals the one-variable series coefficientwise in $z$ and in $L-s_*$. \Cref{lem:gam-deltasrc} on random data with $c\in\{0,\pm 1,2,-3,1/5\}$ through $z^{12}$, the moment table against the Wick recursion up to degree $24$; the wrong sign of $\delta$, a term $m^3$ and a term $s^2m$ break it. For non-integer $M_*=1/5,1/3,-1/2,7/4$, with the Bernoulli-polynomial form of \cref{rmk:gam-Mcomplex}, the series is regular with the shifted value, through $z^{14}$.

\emph{Falsifiers} (\cref{rmk:gam-necessity}~(c)), exactly, $n\leq 12$: $\Gamma(a/z)^2$ gives poles of order exactly $n$ at $z^n$, with leading coefficients $((2n-1)!!)^2(-1/a_*)^n/(n!\,16^n)$ ($\tfrac1{16},\tfrac9{512},\tfrac{75}{8192},\tfrac{3675}{524288}$ at $a_*=-1$); $\Gamma(a/z)\Gamma(-a/z)$ poles of order $2n$; $\Gamma(2a/z)$ and $\Gamma(a/z)\Gamma(1+a/z)$ poles of order $n$; three coincident factors the leading coefficients $\tfrac19,\tfrac2{81},\tfrac{14}{2187},\tfrac8{19683},-\tfrac{440}{177147}$; a perturbation of the coefficient of $(z/a)^{n_0}$ in the degenerate Stirling series a first pole at $z^{n_0}$, of order $n_0$, with coefficient the perturbation divided by $a_*^{n_0}$; an exponent $a^E$, $E\neq M_*-1/2$, a pole at $z^1$ with coefficient $(1/24+\beta_1(M_*)-E^2/2)/a_*$; $\cF=1/(\sigma-s_*)$ poles of order $n$ with coefficient $n!(-1/a_*)^n$ (observed); $\cF(s_*,0)=0$ with $M_*\neq 0$ simple poles; a degenerate argument with a quadratic term coefficients depending on $\log\mathsf Q$; an amplitude depending on $\lambda$ a logarithmic divergence. Positive controls hold: $\cF$ depending holomorphically on $\mathsf Q$; the exponent $+1/2$ with $M_*=0$ (that is $\Gamma(1+w)=w\Gamma(w)$); $(\sigma-s_*)\Gamma(a/z)$.

\emph{The quintic.} In the form of \cref{thm:gam-lmi} at $s_*=1$: no pole and $c(z;1)=h$, exactly through $z^{12}$ and modulo two primes through $z^{16}$, also at the colours $1,2$ with an explicit $\zeta$ modulo a prime, where the value $h(z)$ in place of $h(\zeta^{-i}z)$ fails. With $\Gamma(5\sigma/z)$ in place of $\Gamma(1+5\sigma/z)$ the value at $L=1$ is still $h$, while the $c_n(L)$ change from $n=2$ on (\cref{rmk:gam-blind}); with $\Gamma(1+(1-\sigma)/z)$, that is the measure $\dd x_0$ for the degenerate factor, the series is regular with the value for $\delta=z$, which differs from $h$ at $z^1$; at $L=0$ with $\Gamma(5\sigma/z)$ the value is $\exp(5\sum_{5\mid 2k-1}\bk_kz^{2k-1})\neq\Ginf$; at $L=0$ the Bernoulli polynomials at $0$ and $-1/5$ in place of $1/5$ fail at $z^5$, and at $4/5$ only at $z^{10}$; the twist weight $+5$ in $h$ fails at $z^1$.

\emph{Other targets} (\cref{rmk:gam-general}), exactly: the $\cO(5)$-twisted $\bbP^4$ at four sets of distinct rational weights, all $20$ colours, through $n=24$, with the values $h_{\bw_i}$, $\bw_i=\{\lambda_i-\lambda_j\}_{j\neq i}\cup\{-5\lambda_i\}$; $\bbP^1$ and $\bbP^2$ (two choices of weights each), local $\bbP^2$ (two lifts), $\cO(-1)\oplus\cO(-1)$ over $\bbP^1$ (two lifts), $\cO(3)$, $\cO(4)$, $\cO(6)$ over $\bbP^4$ and $\cO(4)$ with a shifted lift ($12$ theories, $40$ colours, through $n=20$), with the values predicted from the twisted I-functions, in which an Euler-class twist contributes minus its fibre weight and an inverse Euler class plus its fibre weight; the charts at $q=\infty$; random regular data.

\emph{The seven-variable chart} (\cref{rmk:gam-sevenvar}). The chart series equals $c(z;L)$ exactly over $\bbQ(\zeta)$ through $z^3$ at $L=2,-2/3$ and modulo the primes $2^{62}+177$ and $2^{80}+85$ through $z^6$ at $L=3,5/4,-1/2$; at $L=1$ it equals $h$ (exactly through $z^7$); the seven-variable series equals $c(z;L)$ exactly through $z^4$ at $L=2$; without the Jacobian, or with $\dd y/y$, the chart fails at $z^1$, resp.\ $z^2$. Expanded in $q$, the chart reproduces $h$ times the I-function of colour $0$ on all $98$ non-zero coefficients of $q^dz^n$ with $d\leq 6$, $-d\leq n\leq 10$, exactly.

\emph{The formal lemmas of \cref{sec:fsp}} on random data, exactly: \cref{lem:fsp-ibp}, \cref{lem:fsp-cov}, \cref{lem:fsp-fubini} (inner dimensions $1$ and $2$), the formal delta lemma invoked in \cref{rmk:gam-sevenvar} (which is not stated in \cref{sec:fsp}) and \cref{lem:fsp-param} through $z^6$ to $z^8$; \cref{lem:fsp-stirling} for $m=-3,\dots,3$, both measures, at $a=1,-2/7,3/2$, through $z^{40}$ (through $z^{80}$ at $a=1$); \cref{lem:fsp-dlmf5118} as polynomials in $\theta$ through $w^{-60}$, and $B_k(\theta+1)-B_k(\theta)=k\theta^{k-1}$ for $k\leq 61$; the non-integer form of Euler's integral behind \cref{rmk:gam-Mcomplex}, with a symbolic exponent, through $z^{12}$.


